\section{Discussion}
\label{sec:discussion}

\subsection{Lessons}

\paragraph{Varying the compilation.} None of the 58 pipeline findings occurs with a contiguous
float32 tensor on the default path, and 20 of them occur only when the path changes. RQ1 isolates this
effect: with program set, inputs, and oracle held fixed on a corpus already run under three configurations, one
additional path factor produced two defects within minutes. The explanation lies in the code
that each factor reaches. A graph break makes every live local an input of a resume function, so the machinery
that reconstructs Python objects and replays side effects is exercised for objects that a single-graph run
never reconstructs. A dispatch mode routes composite operators through a path on which conjugate views are
materialized lazily. A packaged runtime has a lifetime that the just-in-time runtime does not have. Program
diversity does not reach this code because the program cannot select the path.

\paragraph{Seeds written after a finding.} Both corpus additions of RQ1 followed a finding. The first, written
for the mechanisms behind C41 and C42, re-found those two under the break factor and exposed C44 through its
exception-flow programs on the unmodified path; the second, written for C44, delimited it without adding a defect. A batch of 34 generic programs written before any finding produced nothing new.

\paragraph{Lost validation.} Nine findings return a value where eager execution raises, and 13 raise where
eager execution returns. The first group is not harmless: C37 computes convolutions on Boolean inputs and
softmax on integers with results that no eager kernel defines, and the maintainers opened an umbrella issue on
lost validation during the study. An oracle that ignores exceptions, or a harness that discards every context in
which eager execution raises, misses this class. In the historical benchmark, fourteen open issues ran
zero tests until the harness stopped discarding such seeds.

\paragraph{Reference strength.} Three accuracy findings (C35, C36, C40) are invisible in an
eager-versus-compiled comparison at ordinary inputs and appeared only against a multi-precision reference next to
singularities. A float64 reference separates rounding noise from wrong results for
well-conditioned computations but not for ill-conditioned ones: five closed historical issues that \tool still
flags had been judged intended behavior because a legal re-association is amplified by a later operator.
Such findings are annotated for manual review; condition-aware suppression is future work.

\subsection{Threats to Validity}
\label{sec:threats}

\paragraph{Internal validity.} Harness defects can appear as compiler defects. Two occurred during the study: a
wrapper that compiled the wrong callee, and a sample generator whose \code{no\_grad} finalizer ran inside a
later trace and inserted a grad-mode node into exported graphs. Every candidate is therefore confirmed with a stand-alone script that does not
import \tool, in a fresh interpreter, on two PyTorch versions and two operating systems, and every report is
re-verified from its text before filing. Injected faults may not represent real defects, so the evaluation
also uses mined issues and defects found in current builds.

\paragraph{Construct validity.} A differential oracle detects a divergence and does not decide which side is wrong; ten findings had their root cause outside the pipeline and are counted
separately. Layer attribution was validated against the location of 13 fixes, which biases the check
toward defects that were easy to fix. The maintainer-response counts reflect the tracker at the time of writing;
most path-factor items were filed in the last days before submission. The comparison in RQ5 covers random
generation and NNSmith; TorchProbe, the closest tool, could not be run against our subject.

\paragraph{External validity.} The implementation targets \tcompile; the factor model needs only an eager
reference, pipeline switches, and selectable layers, which other AI compiler stacks also offer but which we have
not evaluated. Most experiments ran on CPUs; the T4 runs confirmed the platform-independent
defects and found CUDA-specific members of C37, but Triton code generation is covered less densely than C++
code generation, and C46 is specific to Windows.
\clearpage
